\lyseis{Β2.2}{Άξονας συμμετρίας σχήματος}
\sol{1} Τέσσερις άξονες: δύο μεσοκάθετοι των πλευρών και οι δύο διαγώνιοι.
\sol{2} α) $2$ \quad β) $2$ \quad γ) $3$ \quad δ) $1$ \quad ε) $0$ \quad στ) $0$.
\sol{3} Πεντάγωνο $5$, εξάγωνο $6$, κύκλος άπειροι.
\sol{4} Δύο: η ευθεία $\mathrm{ΑΒ}$ και η μεσοκάθετος του $\mathrm{ΑΒ}$.
\sol{5} Ο φορέας της διχοτόμου της γωνίας.
\sol{6} Άξονα συμμετρίας έχουν όλα εκτός από τα $\Gamma$ και $\mathrm{Ν}$ (σε μια γραμματοσειρά χωρίς διακοσμητικά).
\sol{7} α) Λάθος. β) Σωστό. γ) Λάθος. δ) Σωστό. ε) Σωστό.
\sol{8} Έξι άξονες, και οι τρεις περνούν από δύο απέναντι κορυφές.
\sol{9} Δύο άξονες: οι μεσοκάθετοι των πλευρών ($1{,}2\neq2{,}1$, άρα δεν είναι τετράγωνο).
\sol{10} Στο κέντρο του ορθογωνίου, όπου τέμνονται οι δύο άξονες.
\sol{11} Ελεύθερη κατασκευή. Π.χ. ένα ορθογώνιο ή έλλειψη.
\sol{12} Δύο άξονες, οι διαγώνιοι του ρόμβου. Τέμνονται στο κέντρο του.
\sol{13} Γιατί όταν διπλώσουμε τη διαγώνιο, οι δύο άλλες γωνίες δεν συμπίπτουν (το τρίγωνο δεν είναι ισοσκελές) και τα δύο τρίγωνα δεν εφαρμόζουν σαν καθρέφτης, αλλά μόνο με περιστροφή.
\sol{14} Δύο: π.χ. το ορθογώνιο (ή ρόμβος). Ακριβώς τρεις: το ισόπλευρο τρίγωνο.
\sol{15} Στο σημείο τομής των αξόνων.
\sol{16} Η δίπλωση αντιστοιχεί την κορυφή στον εαυτό της, άρα ο άξονας διέρχεται από αυτήν και οι δύο ίσες πλευρές εφαρμόζουν· οι γωνίες που σχηματίζουν με τον άξονα είναι ίσες, άρα ο άξονας διχοτομεί τη γωνία.
\sol{17} Ελεύθερη απάντηση.
\sol{18} Ελεύθερη απάντηση.
